\documentclass[11pt]{article}
\usepackage[margin=1in]{geometry}
\usepackage{amsmath,amssymb,amsthm}
\usepackage[T1]{fontenc}
\usepackage{lmodern}
\usepackage[hidelinks]{hyperref}
\usepackage{xcolor}

\newtheorem{theorem}{Theorem}
\newtheorem{corollary}[theorem]{Corollary}

\title{Finite Unions of Weak-Mixing Linear\\Non-Recurrence Sets}
\author{Candidate partial answer to arXiv:1202.3114}
\date{11 August 2026}

\begin{document}
\maketitle

\begin{center}
\fcolorbox{black}{gray!8}{\parbox{0.88\linewidth}{
\textbf{Status: candidate partial result, likely valid.}
This enlarges the known lacunary class but does not settle the arbitrary
lacunary or shifted-rigidity questions.}}
\end{center}

\section{Source question}

Grivaux~\cite{G} asks whether every lacunary sequence $(n_k)$ is a
non-recurrence set for some weakly mixing dynamical system.  The paper proves
this when
\[
 \frac{n_{k+1}}{n_k}\longrightarrow\infty
\]
and, by a different rigidity argument, for divisibility chains
$n_k\mid n_{k+1}$.

We prove that the linear-Gaussian conclusion is stable under finite unions.
This gives a strictly broader ratio condition: it is enough that large jumps
occur within every fixed-length block.

Recall that $D\subset\mathbb N$ is a non-recurrence set for a
measure-preserving system $(X,\mathcal B,\mu,T)$ if there is
$A\in\mathcal B$ with $\mu(A)>0$ and
\[
 \mu(A\cap T^{-n}A)=0\qquad(n\in D).
\]

\section{Finite-union closure}

\begin{theorem}\label{thm:union}
Let $D_1,\ldots,D_r\subset\mathbb N$.  Suppose that for every $j$ there is a
weakly mixing linear dynamical system $(H_j,\mathcal B_j,m_j,T_j)$ with a
non-degenerate invariant Gaussian measure such that $D_j$ is a
non-recurrence set.  Then $D_1\cup\cdots\cup D_r$ is a non-recurrence set for
a weakly mixing linear dynamical system with a non-degenerate invariant
Gaussian measure.
\end{theorem}

\begin{proof}
For each $j$, choose a Borel set $A_j\subset H_j$ of positive measure such
that
\[
 m_j(A_j\cap T_j^{-n}A_j)=0\qquad(n\in D_j).
\]
On the Hilbert direct sum
\[
 H=H_1\oplus\cdots\oplus H_r
\]
consider the block-diagonal bounded operator
$T=T_1\oplus\cdots\oplus T_r$ and the product Gaussian measure
$m=m_1\otimes\cdots\otimes m_r$.  This is again a non-degenerate invariant
Gaussian measure.  A finite product of weakly mixing probability-preserving
systems is weakly mixing, so $T$ is weakly mixing with respect to $m$.

Put $A=A_1\times\cdots\times A_r$.  Then
$m(A)=\prod_jm_j(A_j)>0$.  If $n\in D_1\cup\cdots\cup D_r$, choose $j$ with
$n\in D_j$.  Product measure gives
\[
 m(A\cap T^{-n}A)
 =\prod_{i=1}^r m_i(A_i\cap T_i^{-n}A_i)=0,
\]
because the $j$th factor vanishes.  Thus the union is non-recurrent for the
required weakly mixing linear system.
\end{proof}

The same proof works for arbitrary weakly mixing systems.  The linear and
Gaussian assertions are preserved because only a finite Hilbert direct sum
and a finite product measure are used.

\section{A broader lacunary class}

\begin{corollary}\label{cor:gap}
Let $(n_k)_{k\ge0}$ be a strictly increasing sequence of positive integers.
If there is an integer $r\ge1$ such that
\begin{equation}\label{eq:blockgap}
 \frac{n_{k+r}}{n_k}\longrightarrow\infty,
\end{equation}
then $\{n_k:k\ge0\}$ is a non-recurrence set for a weakly mixing linear
dynamical system with a non-degenerate invariant Gaussian measure.
\end{corollary}

\begin{proof}
For $j=0,\ldots,r-1$, set
\[
 D_j=\{n_{j+\ell r}:\ell\ge0\}.
\]
Along this subsequence, the ratio of consecutive terms is
\[
 \frac{n_{j+(\ell+1)r}}{n_{j+\ell r}},
\]
which tends to infinity by \eqref{eq:blockgap}.  Theorem 1.1 of~\cite{G}
therefore makes each $D_j$ a non-recurrence set for a weakly mixing linear
Gaussian system.  Theorem~\ref{thm:union} applies, and the union of the
$D_j$ is the original set.
\end{proof}

Condition \eqref{eq:blockgap} is strictly weaker than the source's adjacent
ratio hypothesis.  For example, take
\[
 n_{2m}=2^{m^2+m},\qquad n_{2m+1}=2^{m^2+m+1}.
\]
The adjacent ratio $n_{2m+1}/n_{2m}$ is always $2$, so it does not tend to
infinity, whereas $n_{k+2}/n_k\to\infty$.  Corollary~\ref{cor:gap} applies.

More generally, Theorem~\ref{thm:union} combines any finite collection of
classes already known to be weak-mixing linear non-recurrence sets, including
superlacunary streams and divisibility streams.

\section{Why the full question remains}

An arbitrary lacunary sequence with $n_{k+1}/n_k\ge a>1$ need not satisfy
\eqref{eq:blockgap} for any fixed $r$; geometric sequences are the simplest
example.  Some such sequences are covered by divisibility, but a general
finite decomposition into the source's classes is unavailable.

For countably many component sets, the product proof breaks down sharply.
Any positive-measure set disjoint from one of its translates has measure at
most $1/2$, so the product of infinitely many witness measures is zero.
Likewise, the shifted-rigidity argument requires summable rigidity errors in
order to trim all residual intersections; ordinary rigidity gives only
convergence to zero.  These are the two barriers preventing the elementary
closure argument from resolving the source questions in full.

\begin{thebibliography}{9}
\bibitem{G}
S. Grivaux,
\emph{Non-recurrence sets for weakly mixing linear dynamical systems},
Ergodic Theory Dynam. Systems \textbf{34} (2014), 132--152;
arXiv:1202.3114.
\end{thebibliography}

\end{document}
